\documentclass[10pt]{article}
\providecommand{\FIGDIR}{fig}
\newcommand{\DIGESTDATE}{5 October 2026}
\input{preamble}

\begin{document}

% ==================================================================== MASTHEAD
\thispagestyle{empty}
\begin{tikzpicture}[remember picture, overlay]
  \fill[Deep] (current page.north west) rectangle
      ($(current page.north east)+(0,-67mm)$);
  \fill[Amber] ($(current page.north west)+(0,-67mm)$) rectangle
      ($(current page.north east)+(0,-68.4mm)$);
  \foreach \x/\y/\r in {22/12/0.9, 41/26/1.5, 63/9/0.7, 88/31/1.1, 112/17/1.9,
                        138/29/0.8, 159/13/1.3, 178/54/1.0, 196/21/1.6,
                        30/44/0.6, 74/57/1.2, 125/48/0.9, 168/6/0.8, 191/59/0.7,
                        52/18/1.0, 100/60/0.6, 148/20/1.1, 66/49/1.4, 205/45/0.8}
    \fill[Sky, opacity=0.30] ($(current page.north west)+(\x mm,-\y mm)$) circle (\r mm);
\end{tikzpicture}

\vspace*{4mm}
{\sffamily\fontsize{7.6}{9}\selectfont\color{Sky}%
 AUTOMATED DAILY LITERATURE DIGEST\;\textbullet\;PhD RESEARCH WATCH\par}
\vspace{2.2mm}
{\sffamily\bfseries\fontsize{27}{29}\selectfont\color{white}%
 Particulate Matter\par}
{\sffamily\fontsize{27}{29}\selectfont\color{Sky}%
 Monitoring \& Health Effects\par}

\vspace{5.0mm}
{\sffamily\fontsize{8.4}{10}\selectfont\color{white}%
 Issue 2026-10-05\;\;\textcolor{Amber}{\textbar}\;\;Entry window 5 October 2026 (full day)\;\;%
 \textcolor{Amber}{\textbar}\;\;13 records screened in scope\par}
\vspace{18mm}

\noindent
\begin{minipage}[t]{0.190\linewidth}\metric{13}{RECORDS IN SCOPE\\(39 EXCLUDED)}{Deep}\end{minipage}\hfill
\begin{minipage}[t]{0.190\linewidth}\metric{7}{OF 10 SUBTOPIC\\BANDS POPULATED}{Teal}\end{minipage}\hfill
\begin{minipage}[t]{0.190\linewidth}\metric{3}{EFFECT ESTIMATES\\WITH A CI}{Sky}\end{minipage}\hfill
\begin{minipage}[t]{0.190\linewidth}\metric{2}{TIER-A\\RECORDS}{Amber}\end{minipage}\hfill
\begin{minipage}[t]{0.190\linewidth}\metric{5}{RECORDS CARRIED ON\\METADATA ALONE}{Clay}\end{minipage}

\vspace{6mm}

\section*{\sffamily\bfseries\color{Deep}Signal of the day}
\vspace{-1.5mm}{\color{Amber}\rule{\linewidth}{1.1pt}}\vspace{2.5mm}

\begin{keybox}
{\sffamily\bfseries\small\color{Deep}2.1 million adults, 19 million person-years, and the
exposure that carries the dementia signal is the one no low-cost node measures.}\\[1.4mm]
\footnotesize
Batisse and colleagues (\textit{Epidemiology}, tier A) linked high-resolution residential \textbf{ultrafine
particle number} and \textbf{UFP size} to cause-of-death records for 2.1\,million adults in Montreal and
Toronto, 2001--2019: \textbf{20,560} neurodegenerative deaths, and per \textbf{10,000 particles\,cm$^{-3}$}
a dementia-mortality \textbf{HR of 1.22} (1.17--1.27) --- stronger after adjusting for particle size,
positive across every dementia subtype, largest for vascular dementia. \textbf{Caveat:} UFP number and
NO$_2$ share the same traffic gradient, so co-pollutant adjustment cannot cleanly separate them, and
dementia is under-ascribed on death certificates.
\textbf{Why it matters for sensing:} UFP is unregulated, and an optical low-cost sensor is blind to it ---
it reports a scattering proxy for mass, and sub-100\,nm particles scatter almost nothing. A network built
to track PM$_{2.5}$ mass reports approximately zero information about the exposure this cohort implicates.
The same day, a tethered-balloon synthesis over Houston shows the other blind spot: profiles that vary
severely with boundary-layer depth above a surface node that cannot see any of it.\par
\end{keybox}

\vspace{3mm}

\section*{\sffamily\bfseries\color{Deep}What changed today}
\vspace{-1.5mm}{\color{Amber}\rule{\linewidth}{1.1pt}}\vspace{2.5mm}

\footnotesize
\begin{itemize}

\item \textbf{Gene--environment interaction in 313,534 people.} UK Biobank: PM$_{2.5}$ \textbf{HR 1.22 per
5\,$\mu$g\,m$^{-3}$} (1.19--1.26) across five chronic-disease groups, with additive interaction between
high polygenic risk and high PM$_{2.5}$ for neurodegenerative disease (RERI 0.14) and diabetes (0.21).
A healthy-volunteer cohort over a narrow exposure range, five outcomes, four pollutants, and no
multiplicity correction on the interaction terms.

\item \textbf{The 1.5\,m reading is not the 20\,m reading.} A three-level mast inside a Shenyang urban forest
through the heating season found PM$_{2.5}$ and PM$_{10}$ \textit{higher} at 10 and 20\,m than at pedestrian
level, O$_3$ and CO higher at 1.5\,m, NO$_2$ and SO$_2$ peaking mid-canopy. One site, one winter --- but
network siting guidance assumes a well-mixed sub-canopy layer, and this is direct evidence against it.

\item \textbf{Triangulation that is not independent.} Pretoria: hazard quotients above unity in all age
groups, case-crossover \textbf{OR 1.027} (1.006--1.049) per 10\,$\mu$g\,m$^{-3}$, and 158--748 admissions
avertable depending on the counterfactual. The three methods re-use one exposure series, and summing
element-specific attributable cases double-counts the mass those elements are part of.

\item \textbf{Respirable dust is the wrong metric for a silica hazard.} Personal sampling across five
stone-working tasks in Thailand: all means below the Thai OEL and TLV, highest in carving
(0.051\,mg\,m$^{-3}$), modelled ILCR $3.3\times10^{-7}$. No crystalline silica fraction is reported, and an
EPA inhalation-unit-risk framework applied to mixed mineral dust yields a number with no
exposure--response behind it.

\end{itemize}

\vspace{2mm}

\begin{keybox}
{\sffamily\bfseries\footnotesize\color{Deep}Window continuity}\\[1.2mm]
{\fontsize{7.2}{8.8}\selectfont
The 4 October issue closed with \texttt{last\_entry\_date} at 4 October. This issue collects
\textbf{5 October in full}, continuous with it. Harvest and screening were completed on the 6 October run
and cached; this run authored from that cache without re-querying. Dailies for 6--8 October and the W40
weekly (27 September--3 October) remain owed and follow in this backfill.\par}
\end{keybox}

\vspace{3mm}
\par\vskip 0pt plus 18\baselineskip\penalty-200\vskip 0pt plus -18\baselineskip
\section*{\sffamily\bfseries\color{Deep}Corpus analytics}
\vspace{-1.5mm}{\color{Amber}\rule{\linewidth}{1.1pt}}\vspace{2.5mm}

\noindent\includegraphics[width=\linewidth]{\FIGDIR/f1_subtopics.png}
\figcap{Subtopic assignment of the 13 in-scope records. \textbf{Seven of ten bands populated};
\textit{Exposure assessment} leads with \textbf{four} (three of them metadata-only), \textit{Sensing} has
three, \textit{Occupational \& indoor} two, and one each for \textit{Neuro}, \textit{Other clinical},
\textit{Respiratory} and \textit{Burden}. The clinical side is thin but carries both tier-A records.}

\vspace{2mm}
\noindent\includegraphics[width=\linewidth]{\FIGDIR/f2_design_endpoint.png}
\figcap{Left --- architecture: \textbf{five} records carry no abstract and assert no
design; \textbf{four} are measurement campaigns (the canopy mast, the tethered balloon, the Nagpur
monitoring year and the Xinxiang mercury sampling); two are cohorts (CanCHEC, UK Biobank), one acute
(the case-crossover) and one cross-sectional (the stone carvers). Right --- the majority of records report
no health endpoint; the four that do are split between neurodegenerative mortality, multi-system chronic
disease, respiratory admissions and a modelled cancer risk.}

\vspace{2mm}
\noindent\includegraphics[width=\linewidth]{\FIGDIR/f3_metric_geography.png}
\figcap{Left --- particle metric: unspeciated PM and PM$_{2.5}$ + PM$_{10}$ jointly and
PM$_{2.5}$ only take three each, with one record each for size-resolved number distribution (the Houston
profiles), optical properties, settleable/respirable dust (the personal sampling) and bioaerosol. The UFP
cohort is the only record whose exposure is a \textit{number} concentration. Right --- geography: China
\textbf{five}, North America two, and one each for Europe, Sub-Saharan Africa, South Asia, Southeast Asia,
the Middle East \& North Africa and global / multi-region.}

\vspace{2mm}
\noindent\includegraphics[width=\linewidth]{\FIGDIR/f6_lifecourse.png}
\figcap{Life-course windows: older adults one (the UFP dementia-mortality cohort), working age
one (stone carvers under personal sampling), childhood one (the paediatric-ward bioaerosol model, metadata
only). In utero and adolescence have no record. The UK Biobank cohort spans adulthood and is not assigned
to a single window.}

\vspace{2mm}
\noindent\includegraphics[width=\linewidth]{\FIGDIR/f4_heatmap.png}
\figcap{Subtopic $\times$ study architecture for the 13 records. The metadata-only block spans
exposure assessment, sensing and occupational; measurement campaigns occupy sensing, exposure and burden;
the three epidemiological architectures each sit alone in their clinical band.}

\vspace{2mm}

\noindent\includegraphics[width=\linewidth]{\FIGDIR/f5_forest.png}
\figcap{Three estimates with 95\,\% CIs, log scale, from three different cohorts and three different
exposure contrasts --- per 10,000\,particles\,cm$^{-3}$ of UFP number (dementia mortality), per
5\,$\mu$g\,m$^{-3}$ of annual PM$_{2.5}$ (chronic-disease incidence) and per 10\,$\mu$g\,m$^{-3}$ of daily
PM$_{2.5}$ (respiratory admissions). The first two agree at 1.22 by coincidence of scaling, not of effect:
one is a chronic number-concentration contrast and the other an annual-mass contrast, and they are not
poolable. The third is an acute estimate and belongs to a different estimand entirely.}

\vspace{4mm}

\par\vskip 0pt plus 24\baselineskip\penalty-200\vskip 0pt plus -24\baselineskip
\section*{\sffamily\bfseries\color{Deep}Paper-level digest}
\vspace{-1.5mm}{\color{Amber}\rule{\linewidth}{1.1pt}}\vspace{2.5mm}

{\fontsize{7.4}{9}\selectfont\color{Slate}Ordered by cluster. Each entry gives the
methodological core and the finding that matters, not an abstract paraphrase. Records
marked \textit{metadata only} carry no recoverable abstract and assert no finding.\par}

\band{Neuro / mental health\hfill 1 record}{Deep}

\paperentry{Batisse et al. 2026}
{Long-term exposure to ultrafine particles and mortality from neurodegenerative diseases: a population-based cohort study}
{Epidemiology}
{2.1 million adults in the Canadian Census Health and Environment Cohorts, 2001-2019, 19 million person-years and 20,560 neurodegenerative deaths. Per 10,000 particles/cm3 of UFP number, dementia mortality HR 1.22 (95\% CI 1.17-1.27), strengthened after adjustment for UFP size; positive for every dementia subtype, largest for vascular dementia. Models were run with and without inverse-probability censoring weights for competing events and adjusted for co-pollutants. Limitation: postal-code exposure from a land-use-regression UFP surface, no mobility or indoor time, and cause of death on the death certificate under-ascribes dementia; UFP number and NO$_2$ share the traffic gradient, so co-pollutant adjustment cannot fully separate them. Relevance: high - UFP is unregulated and unmeasured by every optical low-cost node, which sees mass, not number; this is the strongest current argument for CPC-class instruments in dense networks.}
{10.1097/ede.0000000000002054}

\par\vskip 0pt plus 10\baselineskip\penalty-200\vskip 0pt plus -10\baselineskip
\band{Respiratory \& allergic\hfill 1 record}{Sky}

\paperentry{Howlett-Downing and Wichmann 2026}
{Triangulating health risks of PM$_{2.5}$ and trace elements in Pretoria: evidence from health risk assessment, case-crossover, and concentration-response function analyses}
{Int J Environ Health Res}
{Three methods applied to the same secondary dataset: hazard quotients, a case-crossover analysis and concentration-response functions. Hazard quotients exceeded unity in adults, children and infants; the case-crossover OR was 1.027 (95\% CI 1.006-1.049) per 10 $\mu$g\,m$^{-3}$ PM$_{2.5}$. Meeting the South African standard would avert 158 admissions annually, rising to 342 and 748 under the WHO daily and annual guidelines; element-specific functions give 110-608. Limitation: the three methods are not independent evidence - they re-use one exposure series from a small number of sampling days, and summing element-specific attributable cases double-counts the mass they are part of; hazard quotients are screening-level and not comparable to epidemiological risk. Relevance: moderate - Southern African monitoring is thin, and a saturation low-cost network would resolve the within-city contrast this design cannot.}
{10.1080/09603123.2026.2733112}

\par\vskip 0pt plus 10\baselineskip\penalty-200\vskip 0pt plus -10\baselineskip
\band{Exposure assessment \& modelling\hfill 4 records}{Clay}

\paperentry{Chavan and Lataye 2026}
{Application of Openair package in R programming and PCA: unveiling air quality dynamics and meteorological influences}
{Environ Sci Pollut Res Int}
{A full year of regulatory monitoring data for a tier-2 Indian city analysed with openair and PCA. Winter AQI reached 492 (severe) in February with monsoon improvement; varimax PCA gave two components explaining 73.1\% of variance, read as anthropogenic combustion (PM, NO$_2$) and photochemistry (O$_3$, temperature); PM$_{2.5}$ and PM$_{10}$ exceeded CPCB standards on 27-29 days per winter month. Limitation: descriptive reanalysis of one station-year with no source apportionment, no receptor model and no validation - PCA components here are correlation structure, not sources, and the policy recommendations do not follow from the analysis. Relevance: low as method, moderate as a reminder that tier-2 cities have one regulatory station and no spatial information at all.}
{10.1007/s11356-026-38275-w}

\paperentry{Zhao et al. 2026c}
{Spatial Heterogeneity and Multi-scale Statistical Associations of Urban Air Pollution Revealed by Interpretable Machine Learning in Shenzhen, China}
{Atmos Pollut Res}
{\textsc{metadata only.} Metadata-only record: title, authors and DOI deposited 5 Oct with no abstract. By title, an interpretable machine-learning treatment of intra-urban spatial heterogeneity at several scales in Shenzhen. Scale dependence of the predictor set is the question that decides how dense a network has to be before adding nodes stops buying information, so the result matters more than the method. Relevance: potentially high; on the retry list.}
{10.1016/j.apr.2026.103230}

\paperentry{Hopke et al. 2026}
{Characterization of PM Pollution Episodes in Kuwait}
{Atmos Pollut Res}
{\textsc{metadata only.} Metadata-only record: title, authors and DOI deposited 5 Oct with no abstract. By title, a characterisation of PM pollution episodes in Kuwait. Dust-dominated Gulf environments are where optical low-cost sensors fail hardest - coarse mineral dust sits outside the calibrated size and refractive-index range of nephelometric nodes - and the episode statistics are the input any regional network design needs. Relevance: potentially high; on the retry list.}
{10.1016/j.apr.2026.103226}

\paperentry{Shi et al. 2026}
{Compositional transition and coupled formation pathways of PM$_{2.5}$ pollution in a typical industrial city: An empirical study based on Handan, China}
{Environ Pollut}
{\textsc{metadata only.} Metadata-only record: Elsevier deposited title, authors and DOI on 5 Oct with no abstract. By title, an empirical account of how PM$_{2.5}$ composition in a heavily industrial Hebei city has shifted and how its formation pathways couple. A compositional transition from primary to secondary is the slow change that silently alters the relationship between an optical sensor's scattering signal and gravimetric mass over years of a network's life. Relevance: moderate; on the retry list.}
{10.1016/j.envpol.2026.129296}

\par\vskip 0pt plus 10\baselineskip\penalty-200\vskip 0pt plus -10\baselineskip
\band{Sensing, forecasting \& instrumentation\hfill 3 records}{Amber}

\paperentry{Chang et al. 2026}
{Wintertime vertical stratification of air pollutants and meteorological associations within an urban forest canopy in Shenyang, China}
{Environ Pollut}
{A three-level monitoring mast inside an urban forest in a cold-climate northern Chinese city through the winter heating season. PM$_{2.5}$ and PM$_{10}$ were generally higher at 10 and 20 m than at pedestrian level, O$_3$ and CO higher at 1.5 m, and NO$_2$ and SO$_2$ peaked mid-canopy; relative humidity and wind speed dominated the statistical and machine-learning associations, but height-dependently. Limitation: one site, one winter, no replication mast and no flux measurement, so the inversion of the usual near-ground maximum is described rather than explained - and at wintertime humidity, uncorrected optical PM readings are themselves height-correlated through the humidity profile. Relevance: high - siting guidance for low-cost networks assumes a well-mixed layer below canopy height; this is direct evidence that the 1.5 m reading is not the 20 m reading.}
{10.1016/j.envpol.2026.129270}

\paperentry{Mei et al. 2026}
{Synthesis of the tethered balloon system and other TRACER campaign measurements elucidates aerosol property profiles}
{Atmos Chem Phys}
{149 tethered-balloon flights over greater Houston in summer during the DOE ARM TRACER campaign, with back-trajectory k-means clustering into marine, mixed and urban/long-range air masses. Profiles vary strongly with cluster, modulated by boundary-layer depth and coastal circulation; the marine cluster had CCN at 0.8\% supersaturation below 1,000 cm$^{-3}$, urban and mixed clusters higher. Limitation: one coastal city, one summer, and CCN is inferred from size distribution under an assumed kappa rather than measured, so composition heterogeneity propagates straight into the activated fraction; tethered-balloon sampling is intermittent and fair-weather-biased. Relevance: high - this is the measurement that exposes the central fiction of a surface network, namely that a 2 m reading represents the column a satellite retrieval sees.}
{10.5194/acp-26-13885-2026}

\paperentry{Hazelwood et al. 2026}
{Size-Resolved Particle Dry Deposition Velocities in a Controlled Spherical Chamber: Effects of Turbulence, Surface Complexity, and Particle Physicochemical Properties}
{J Aerosol Sci}
{\textsc{metadata only.} Metadata-only record: Elsevier deposited title, authors and DOI on 5 Oct with no abstract. By title, a chamber determination of size-resolved dry deposition velocity against turbulence intensity, surface complexity and particle physicochemical properties. Deposition velocity is the term that sets indoor particle loss rates and therefore the inferred air-change rate in every filtration study that does not measure one, and it is the largest uncertainty in indoor-outdoor infiltration models. Relevance: potentially high; on the retry list.}
{10.1016/j.jaerosci.2026.106903}

\par\vskip 0pt plus 10\baselineskip\penalty-200\vskip 0pt plus -10\baselineskip
\band{Occupational \& indoor\hfill 2 records}{Amber}

\paperentry{Ketsakorn and Chaiyadej 2026}
{Probabilistic health risk assessment of respirable dust exposure among stone carvers in Thailand using average daily dose and Monte Carlo simulation}
{PLoS One}
{Personal sampling across mining, cutting, splitting, knocking and carving, with a 10,000-iteration Monte Carlo ILCR using US EPA exposure parameters. Mean respirable dust was below both the Thai OEL and the TLV everywhere, highest in carving at 0.051 mg/m3; ILCR stayed under 1e-6, highest in carving at 3.3e-7. Limitation: respirable dust mass is the wrong metric for a silica hazard - no crystalline silica fraction is reported, and an EPA inhalation-unit-risk framework applied to a mixed mineral dust produces a number with no established exposure-response behind it; small sample, one season. Relevance: low-moderate - the task-level contrast is exactly what a wearable logging monitor resolves and a single shift-average sample does not.}
{10.1371/journal.pone.0359706}

\paperentry{Zhao et al. 2026d}
{Behavior-resolved bioaerosol exposure modeling reveals scale-dependent occupational risk in pediatric wards}
{Build Environ}
{\textsc{metadata only.} Metadata-only record: Elsevier deposited title, authors and DOI on 5 Oct with no abstract. By title, a bioaerosol exposure model in paediatric wards resolved by occupant behaviour, reporting scale-dependent occupational risk. Behaviour-resolved models are the class of work that decides whether a fixed room monitor can stand in for a worker's exposure at all, and the paediatric ward is a setting where that substitution is routinely assumed. Relevance: potentially high; on the retry list.}
{10.1016/j.buildenv.2026.115364}

\par\vskip 0pt plus 10\baselineskip\penalty-200\vskip 0pt plus -10\baselineskip
\band{Burden, policy \& mitigation\hfill 1 record}{Coral}

\paperentry{Zhao et al. 2026b}
{Pollution Characteristics, Influencing Factors and Health Risk Assessments of Mercury in Atmospheric PM$_{2.5}$ from Xinxiang, China}
{Atmosphere}
{Year-long filter sampling gives an annual mean PM$_{2.5}$-bound Hg of 0.399 ng/m3, significantly higher in winter (p<0.001) and in the heating than the non-heating period, positively correlated with SO$_2$, CO, PM$_{2.5}$, OC, EC, Cl-, Zn and Ca and negatively with temperature, O$_3$ and V - a coal-combustion signature. Hazard index 2.06e-3 (children) and 8.46e-4 (adults), both far below 1, with ingestion the dominant modelled pathway. Limitation: the samples are from 2015 and describe an emissions regime two major control phases out of date; an ingestion-dominant pathway for an airborne metal is an artefact of the EPA soil-oriented default parameters. Relevance: low - logged as a composition time point, not as a risk estimate.}
{10.3390/atmos17100974}

\par\vskip 0pt plus 10\baselineskip\penalty-200\vskip 0pt plus -10\baselineskip
\band{Other clinical endpoints\hfill 1 record}{Slate}

\paperentry{Tian et al. 2026}
{Annual average exposure to ambient air pollution, polygenic risk, and incidence of major chronic diseases: a large-scale prospective cohort study}
{Hum Genet}
{313,534 UK Biobank participants with LUR annual pollutant estimates and weighted polygenic risk scores for each of five disease groups; Cox models with RERI and AP for additive interaction. PM$_{2.5}$ carried the strongest association, HR 1.22 per 5 $\mu$g\,m$^{-3}$ (95\% CI 1.19-1.26); high genetic risk acted independently (diabetes HR 1.77 per SD of PRS), with significant additive interaction between high PRS and high PM$_{2.5}$ for neurodegenerative disease (RERI 0.14) and diabetes (RERI 0.21). Limitation: UK Biobank is a healthy-volunteer cohort with a narrow PM$_{2.5}$ range, the LUR surface is a single baseline-era estimate carried forward, and five outcomes with four pollutants invites selective emphasis - the interaction terms are modest and not corrected for multiplicity. Relevance: moderate - gene-environment interaction raises the value of resolving exposure well at the individual level, which is the case for personal rather than ambient monitoring.}
{10.1007/s00439-026-02875-w}

\clearpage
\section*{\sffamily\bfseries\color{Deep}Corpus register}
\vspace{-1.5mm}{\color{Amber}\rule{\linewidth}{1.1pt}}\vspace{2.5mm}

\input{table_rows}

\vspace{2mm}

\begin{keybox}
{\sffamily\bfseries\footnotesize\color{Deep}Provenance and method}\\[1.2mm]
{\fontsize{7.2}{8.8}\selectfont
Entry window \textbf{5 October 2026}, single day, continuous with the 4 October issue. Harvest
and screening for this day were completed on the \textbf{6 October} run and cached
(\texttt{cache/2026-10-05/\_fresh.json}); this run authored from that cache without re-querying.
\textbf{52} fresh candidates after cross-day deduplication, \textbf{13 admitted}, \textbf{39 rejected} with
logged reasons (\textit{ACP} cloud microphysics and radiation, ozone-only and NO$_2$-only analyses, water
and soil chemistry, life-cycle assessment, and the usual PubMed broad-axis noise). \textbf{Five records
carried on metadata alone} (38\,\%, a series high) --- Elsevier's 5 October deposit is abstract-free across
four journals, and the retry leg recovered none. \textbf{Consensus} returned 10 calibration hits,
\textbf{0 new}. \textbf{Scholar Gateway} fails with \texttt{ACCESS\_DENIED}; it has now been unavailable
for the whole of this backfill. \textbf{ClinicalTrials.gov}: no new PM or air-pollution registration in the
window. DOI gate: \textbf{0 fail, 0 warn}. Abstracts remain the copyright of their respective
publishers.\par}
\end{keybox}

\end{document}
